\documentclass[10pt,a4paper]{amsart}
\usepackage[margin=19mm]{geometry}
\usepackage{amsmath,amssymb,mathtools}
\usepackage[scr=rsfs,cal=euler]{mathalfa}
\usepackage{xfrac}
\usepackage[unicode,hidelinks,bookmarks=false]{hyperref}
\newcommand{\ie}{i.e.\ }
\newcommand{\cf}{cf.\ }
\newcommand{\resp}{respectively\ }
\newcommand{\Subsection}{\subsection}
\providecommand{\nobreakdash}{\nobreak\hbox{-}}
\setlength{\emergencystretch}{2em}
\multlinegap0pt
\usepackage{bm}
\usepackage{bbm}
\usepackage{dsfont}
\usepackage{xspace}
\usepackage{cancel}
\usepackage{mathtools}
\usepackage{amsthm}
\makeatletter
\@ifundefined{lemma}{\newtheorem{lemma}{Lemma}}{}
\makeatother
\DeclareMathOperator*{\argmin}{arg\,min}
\newcommand{\B}{\mathcal{B}}
\newcommand{\dborel}{\mathcal{B}_{\dspace}}
\newcommand{\dspace}{\mathcal{D}}
\newcommand{\initmeasomega}{P_{\text{init}, \omega}}
\newcommand{\initmeas}{P_{\text{init}}}
\newcommand{\obsmeas}{P_{\text{obs}}}
\newcommand{\pborel}{\mathcal{B}_{\pspace}}
\newcommand{\predictmeas}{P_\text{pred}}
\newcommand{\pspace}{\Lambda}
\newcommand{\p}{\mathbb{P}}
\newcommand{\updatemeas}{P_{\text{up}}}
\title[Iterative Data-Consistent Inversion with Multiple Push-forwa]{Iterative Data-Consistent Inversion with Multiple Push-forward Constraints}
\author{Tianyi Jiang, Troy Butler, Timothy Wildey, Tim Kutta, Haonan Wang}
\date{Source: arXiv:2603.00033v2 (CC BY 4.0)}
\begin{document}
\maketitle

\noindent\emph{Source and licence.} Iterative Data-Consistent Inversion with Multiple Push-forward Constraints. Version arXiv:2603.00033v2 (CC BY 4.0), distributed under the Creative Commons Attribution 4.0 International licence, as recorded in the arXiv licence metadata for this exact version and shown on the abstract page https://arxiv.org/abs/2603.00033v2. Full rights notice: licenses/arxiv-2603.00033-CC-BY-4.0.txt.\par\smallskip

\noindent\textit{Editorial scope.} Counted unit: the Lemma whose source label is \texttt{lemma:for\_bac} in the article (source0.tex lines 297--308, source label \texttt{lemma:for\_bac}), delivered together with the article's own complete original proof of it (source0.tex lines 311--396, one \texttt{proof} environment of 86 lines). No mathematics is written, repaired or re-derived: statement and proof are the article's own text line for line. Required context retained uncounted: integration (lines 221--223). Every definition, display and label of those lines is retained unchanged. Omitted from this package and not redistributed: 1--220, 224--296, 309--310, 397--1658. Nothing inside a retained excerpt is omitted or altered: every retained body is delivered exactly as the article's lines are written, so no omission receipt of any kind accompanies this package. Citations: the retained excerpts carry no citation command at all, so no bibliography entry of the article is transcribed and no reference is redistributed. Reference adaptation: none was needed and none was made; every reference inside the delivered unit resolves inside it, so no unresolved reference or citation is produced. Printed slips kept literal: no printed word, symbol, hypothesis or number of the counted statement or of its proof was altered. Layout and class: the article's own class is \texttt{article} with packages including lineno, verbatim, amsmath, amsfonts, epsfig, graphics, setspace, amssymb, amsbsy, float, stmaryrd, multirow, caption, enumerate; the wrapper is the lane's pinned \texttt{amsart} editorial wrapper at 10pt on A4, which restates the theorem-like environments the retained text needs and adds the article's own macro definitions that the retained text needs (\texttt{\textbackslash B}, \texttt{\textbackslash argmin}, \texttt{\textbackslash dborel}, \texttt{\textbackslash dspace}, \texttt{\textbackslash initmeas}, \texttt{\textbackslash initmeasomega}, \texttt{\textbackslash obsmeas}, \texttt{\textbackslash p}, \texttt{\textbackslash pborel}, \texttt{\textbackslash predictmeas}, \texttt{\textbackslash pspace}, \texttt{\textbackslash updatemeas}), with extra wrapper packages (bm, bbm, dsfont, xspace, cancel, mathtools). The delivered numbering is the unit's own. Assets: no retained span contains a figure environment, a graphics-inclusion command, a PDF-page inclusion, a table or a TikZ picture; no image file is copied into this package, the package declares no asset dependency and no image file is redistributed with it. Licence: version arXiv:2603.00033v2 of this article is distributed under the Creative Commons Attribution 4.0 International licence, as recorded in the arXiv licence metadata for this exact version and shown on the abstract page https://arxiv.org/abs/2603.00033v2. A full rights notice is delivered at licenses/arxiv-2603.00033-CC-BY-4.0.txt. Characters: every retained excerpt is pure ASCII, so the delivered engine, XeLaTeX with its default Latin Modern text font, reports no missing character.
\begin{equation} \label{integration}
    P_\text{up}(\cdot) := \int_{\mathcal{D}} P_{\text{init}, \omega}(\cdot) dP_\text{obs}(\omega).   
\end{equation}

\begin{lemma}\label{lemma:for_bac} 
Let $\obsmeas$ denote a probability measure on $(\dspace, \dborel)$ and $\initmeas$ denote any measure on $(\pspace,\pborel)$ such that $\obsmeas$ is absolutely continuous with respect to $\predictmeas=\initmeas\circ \phi^{-1}$. 
Let $\p = \{P\ll \initmeas\, : \, P\circ \phi^{-1} = \obsmeas\}$, then $\updatemeas\in \p$ as constructed in~\eqref{integration} solves both the forward minimization and backward minimization problems, i.e., 
\[
    \updatemeas = \argmin_{P\in \p} D_{f}(P || \initmeas) \quad \text{Forward minimization},
\]
and, if $\initmeas \ll \updatemeas$, then
\[
    \updatemeas = \argmin_{P\in \tilde{\p}} KL(\initmeas || P) \quad \text{Backward minimization}
\]
where $\tilde{\p}$ is the subspace of $\p$ such that $\initmeas \ll P$ for all $P\in\tilde{\p}$.
\end{lemma}

\begin{proof}
Note that $\updatemeas\in\p$ is immediate from its construction. 
We therefore first prove that $\updatemeas$ is the solution to the forward minimization problem.
First, we establish that $\updatemeas$ is absolutely continuous with respect to $\initmeas$ so that the standard f-divergence can be computed. 
By the construction of $\updatemeas$ via the disintegration theorem and the Radon-Nikodym formula for applying change-of-measure, we have for each $A\in \B_\Lambda$ the following, 
\begin{eqnarray*}
    \updatemeas(A) =\int_A \, d\updatemeas(\lambda) &=& \int_\dspace \initmeasomega(A) \, d\obsmeas(\omega)\\
                &=&\int_{\dspace} \initmeasomega(A) \cfrac{d\obsmeas}{d\predictmeas}(\omega)\, d\predictmeas(\omega) \\
                &=&\int_A  \cfrac{d\obsmeas}{d\predictmeas}(\phi(\lambda)) dP_\text{init}.
\end{eqnarray*}
It follows that for a.e.~$\lambda$, we have
\[
    \cfrac{d\updatemeas}{d\initmeas}(\lambda) = \cfrac{d\obsmeas}{d\predictmeas}(\phi(\lambda)).
\]
In other words, the Radon-Nikodym derivative of $\updatemeas$ with respect to $P_\text{init}$ can be computed in terms of the Radon-Nikodym derivative of $\obsmeas$ with respect to $\predictmeas$, which exists by the assumption of absolute continuity of $\obsmeas$ with respect to $\predictmeas$. 
Then, by the definition of f-divergence for the non-singular case and the disintegration theorem, we have
\begin{align}\label{eq:Aint}
    D_{f}(\updatemeas || \initmeas) &=  \int_{\Lambda} f\left(\cfrac{d\updatemeas}{d\initmeas}(\lambda)\right) d\initmeas(\lambda) \nonumber \\
    &= \int_{\mathcal{D}} \underbrace{\int_{\phi^{-1}(\omega)} f\left(\cfrac{d\updatemeas}{d\initmeas}(\lambda)\right) d\initmeasomega(\lambda)}_{=: a(\omega)} d\predictmeas(\omega) \\ 
    &= \int_{\mathcal{D}} a(\omega) d\predictmeas(\omega) \nonumber.
\end{align}

Of particular note is that since $\updatemeas$ and $P_\text{init}$ admit the same family of disintegrated conditional densities, then for each $\omega$, the Radon-Nikodym derivative $\frac{d\updatemeas}{d\initmeas}$ is a constant for all $\lambda\in\phi^{-1}(\omega)$.
Thus,
\begin{eqnarray*}
    a(\omega) &=& \int_{\phi^{-1}(\omega)} f\left(\cfrac{d\updatemeas}{d\initmeas}(\lambda)\right) d\initmeasomega(\lambda) \\
            &=& \int_{\phi^{-1}(\omega)} f\left(\cfrac{d\obsmeas}{d\predictmeas}(\phi(\lambda))\right) d\initmeasomega(\lambda)  \\
            &=& f\left(\cfrac{d\obsmeas}{d\predictmeas}(\omega)\right) \underbrace{\int_{\phi^{-1}(\omega)} d\initmeasomega(\lambda)}_{=1}\\
            &=& f\left(\cfrac{d\obsmeas}{d\predictmeas}(\omega)\right).
\end{eqnarray*}

Consider any other $P\in\p$ where it similarly follows that
\begin{align}\label{eq:Bint}
  D_{f}(P || \initmeas) =  \int_{\Lambda} f\left(\cfrac{dP}{d\initmeas}\right) d\initmeas(\lambda) &= \int_{\mathcal{D}} \underbrace{\int_{\phi^{-1}(\omega)} f\left(\cfrac{dP}{d\initmeas}\right) d\initmeasomega(\lambda)}_{=: b(\omega)} d\predictmeas(\omega) \\
  &= \int_{\mathcal{D}} b(\omega) d\predictmeas(\omega). \nonumber 
\end{align}
Due to the convexity of $f$,
\begin{equation*}
    b(\omega) = \int_{\phi^{-1}(\omega)} f\left(\cfrac{dP}{d\initmeas}(\lambda)\right) d\initmeasomega(\lambda) \ge f\left(\int_{\phi^{-1}(\omega)} \cfrac{dP}{d\initmeas}(\lambda) d\initmeasomega(\lambda)\right).
\end{equation*}
Note that for any $B \in \mathcal{B}_{\mathcal{D}}$, applying the disintegration theorem and using the Radon-Nikodym derivative of $P$ with respect to $P_\text{init}$ to perform a change of measure on $\phi^{-1}(B)\in\B_\Lambda$ yields
\begin{equation*}
    P(\phi^{-1}(B)) = \int_{\phi^{-1}(B)} \cfrac{dP}{d\initmeas}(\lambda) d\initmeas(\lambda) = \int_{B} \left( \int_{\phi^{-1}(\omega)} \cfrac{dP}{d\initmeas}(\lambda) \, d\initmeasomega(\lambda) \right) \, d\predictmeas(\omega).
\end{equation*}
At the same time, since $P\in\p$, we have
\begin{equation*}
    P(\phi^{-1}(B)) = \obsmeas(B) = \int_B \, d\obsmeas(\omega) = \int_{B} \cfrac{d\obsmeas}{d\predictmeas}(\omega) d\predictmeas(\omega).
\end{equation*}
Combining these results, it follows that for a.e.~$\omega$, 
\begin{equation*}
    \int_{\phi^{-1}(\omega)} \cfrac{dP}{d\initmeas}(\lambda) d\initmeasomega(\lambda) = \cfrac{d\obsmeas}{d\predictmeas}(\omega).
\end{equation*}
In other words, the integration of $\cfrac{dP}{dP_\text{init}}(\lambda)$ with respect to the conditional measure $\initmeasomega$ over $\phi^{-1}(\omega)$ is equal to $\cfrac{d\obsmeas}{d\predictmeas}(\omega)$.
This is different than the relationship discussed above between $\cfrac{d\updatemeas}{dP_\text{init}}(\lambda)$ and $\cfrac{d\obsmeas}{d\predictmeas}(\omega)$ since $\cfrac{dP}{dP_\text{init}}(\lambda)$ can vary across the set $\lambda\in\phi^{-1}(\omega)$.  
More importantly, this establishes the desired inequality of $b(\omega) \geq f\left(\cfrac{d\obsmeas}{d\predictmeas}(\omega)\right)$.
We have thus demonstrated $\omega$-a.e. that $a(\omega) \le b(\omega)$ so  that
\begin{equation*}
 D_{f}(\updatemeas || \initmeas) = \int_{\mathcal{D}} a(\omega) d\predictmeas(\omega)  \le \int_{\mathcal{D}} b(\omega) d\predictmeas(\omega) = D_{f}(P || \initmeas) ,
\end{equation*}
proving the assertion of forward minimization. 

We next show that $\updatemeas$ is the solution to the backward minimization problem where we restrict $P\in\tilde{\p}$.
We begin by manipulating Radon-Nikodym derivatives through the chain rule to obtain
\begin{align*}
KL( \initmeas || P) - KL( \initmeas || \updatemeas) &= \int_{\Lambda} \log\left(\cfrac{d\initmeas}{dP}(\lambda)\right) d\initmeas(\lambda) - \int_{\Lambda} \log\left(\cfrac{d\initmeas}{d\updatemeas}\right)(\lambda) d\initmeas(\lambda)\\
&= \int_{\Lambda} \log\left(\cfrac{d\updatemeas}{dP}(\lambda)\right) \, d\initmeas(\lambda).
\end{align*}
For each $A\in\B_\Lambda$, define the measure $P^{\prime}$ as follows %{\bf to-do: be more precise about the use of RN derivatives below by referring to the assumptions of absolute continuity that need to be made explicit in this lemma that allow for $dP/d\updatemeas$ to have meaning}
\begin{equation}\label{eq:change_of_conditionals}
    P^{\prime}(A) := \int_A \cfrac{dP}{d\updatemeas}(\lambda) \, d\initmeas(\lambda).
\end{equation}
By direct substitution, 
\begin{eqnarray*}
    KL( \initmeas || P) - KL( \initmeas || \updatemeas) &=& \int_{\Lambda} \log\left(\cfrac{d\updatemeas}{dP}(\lambda)\right) d\initmeas(\lambda) \\
    &=& \int_{\Lambda} \log\left(\cfrac{d\initmeas}{dP^{\prime}}(\lambda)\right) d\initmeas(\lambda) = KL(P_\text{init} || P^{\prime}).
\end{eqnarray*}    
If $P^{\prime}$ is a probability measure, then $KL(P_\text{init} || P^{\prime})\geq 0$ by properties of the KL divergence, which will complete the proof. 
Clearly, $\int \cfrac{dP}{d\updatemeas} d\initmeas$ defines a measure, so we only show that $P^{\prime}(\Lambda) = 1$, which follows from 
\begin{align*}
    P^{\prime}(\Lambda) = \int_{\Lambda} \cfrac{dP}{d\updatemeas}(\lambda) d\initmeas(\lambda) &= \int_{\mathcal{D}} \left(\int_{\phi^{-1}(\omega)} \cfrac{dP}{d\updatemeas}(\lambda) \, d\initmeasomega(\lambda)\right) \, d\predictmeas(\omega) \\
    &= \int_{\mathcal{D}} \left(\int_{\phi^{-1}(\omega)} \cfrac{dP_\omega}{d\initmeasomega}(\lambda) \, d\initmeasomega(\lambda) \right)\, d\predictmeas(\omega) \\
    &= \int_{\mathcal{D}} \int_{\phi^{-1}(\omega)} \, dP_\omega(\lambda) \, d\predictmeas(\omega) \\
    &= P(\pspace) = 1.
\end{align*}
Above, the conditionals defining $dP_\omega$ are uniquely determined by the disintegration of $P$ and we utilized the fact that, by construction, $\updatemeas$ has the same conditionals as $\initmeas$. 
\end{proof}

\end{document}
